% Generated draft: requires mathematical and source review.
\documentclass[11pt]{article}
\usepackage[margin=25mm]{geometry}
\usepackage{fontspec}
\setmainfont{Latin Modern Roman}
\usepackage{amsmath,amssymb,mathtools,unicode-math}
\setmathfont{Latin Modern Math}
\usepackage{enumitem,needspace,xurl,hyperref}
\setlength{\parindent}{0pt}
\setlength{\parskip}{5pt}
\setlength{\emergencystretch}{3em}
\newcommand{\N}{\mathbb{N}}
\newcommand{\Z}{\mathbb{Z}}
\newcommand{\Q}{\mathbb{Q}}
\newcommand{\R}{\mathbb{R}}
\newcommand{\C}{\mathbb{C}}
\newcommand{\F}{\mathbb{F}}
\newcommand{\A}{\mathbb{A}}
\newcommand{\PP}{\mathbb{P}}
\newcommand{\E}{\mathbb{E}}
\newcommand{\eps}{\varepsilon}
\newcommand{\dd}{\,\mathrm{d}}
\begin{document}
\section*{UnsolvedMath ID: 20001976}
\textbf{Source problem code:} catalogue entry
\subsection{Definitions and mathematical statement}
R. Goldin's source asks, for smooth toric varieties \(\C^n//T\), whether ``it'' holds for all toric varieties. The antecedent is absent. The supplied aid interprets the intended issue as integral Kirwan surjectivity with an additive splitting. This interpretation is retained conditionally; it does not recover the original unspecified property or prove a statement for all singular or noncompact toric varieties.

Let a compact connected torus \(T\) act linearly and Hamiltonianly on \(\C^n\), with a fixed inner product on \(\mathfrak t^*\), and let its shifted moment map be
\[
 \mu(z)=A(|z_1|^2/2,\ldots,|z_n|^2/2)-\alpha,
 \qquad A:\R^n\to\mathfrak t^*.
\]
Assume \(\mu^{-1}(0)\) is nonempty and compact. The aid states
\[
 \ker A\cap\R_{\geq0}^n=\{0\},\qquad
 \mu\text{ is proper},\qquad
 \operatorname{Crit}(\|\mu\|^2)\text{ is compact}.
\]
For clarity, a nonnegative kernel direction would give an unbounded ray in the moment level's squared-radius coordinates. Absence of such directions bounds \(\|Av\|\) below by a positive multiple of \(\|v\|\) on the positive cone. At a norm-square critical point with coordinate support \(I\), \(\mu\) is orthogonal to \(\operatorname{im}A_I\) and satisfies
\(\mu=-\operatorname{proj}_{(\operatorname{im}A_I)^\perp}\alpha\).
There are only finitely many supports, and properness then bounds the critical set.

If zero is regular and the action on the level is free, the compact smooth quotient \(X=\mu^{-1}(0)/T\) has Kirwan map
\[
 \kappa:H_T^*(\C^n;\Z)\longrightarrow H^*(X;\Z).
\]
The aid reports that the applicable noncompact extension of Pomerleano--Teleman gives a surjection with a degreewise additive right inverse. This is not a ring splitting. It reports an analogous statement with finite stabilizers after inverting their least common multiple; the precise stabilizer family and theorem hypotheses are not supplied, so no unrestricted integral orbifold claim is made.

Independently, in an effective compact Delzant presentation let \(N\) be the residual torus lattice and
\[
 0\longrightarrow\mathfrak t_{\Z}\longrightarrow\Z^n
 \xrightarrow{\beta}N\longrightarrow0,
 \qquad\beta(e_i)=\nu_i,
\]
where \(\nu_i\) are primitive normals to the polytope inequalities, and \(N^*=\operatorname{Hom}(N,\Z)\). Give \(x_i\) cohomological degree two and set
\[
 J=\left(\sum_i\langle m,\nu_i\rangle x_i:m\in N^*\right),
 \qquad
 I_\Delta=\left(\prod_{i\in S}x_i:
                   \bigcap_{i\in S}F_i=\varnothing\right),
\]
where \(F_i\) is the corresponding equality face. The aid identifies
\[
 H_T^*(\C^n;\Z)=\Z[x_1,\ldots,x_n]/J,\qquad
 H^*(X;\Z)=\Z[x_1,\ldots,x_n]/(J+I_\Delta),
\]
and \(\kappa\) is the quotient map taking coordinate characters to divisor classes. This gives a direct integral surjectivity proof in this smooth compact toric setting. If the presentation includes a strictly redundant inequality, its equality face is empty and \(x_i\in I_\Delta\), so that coordinate class vanishes rather than being counted as a genuine divisor. The quotient's free even-degree integral cohomology also permits an additive splitting; no canonical or multiplicative section is asserted.

Compactness of the nonempty level and freeness/Delzant regularity are essential hypotheses of these conclusions. The source's phrase ``all toric varieties'' leaves open whether singular, noncomplete, orbifold or other presentations are intended, and its missing antecedent prevents a complete general formulation.
\subsection{Short English statement}
For smooth compact free toric reductions, the aid gives integral Kirwan surjectivity and an additive splitting, with an explicit divisor-ring quotient. Compactness of one nonempty linear moment level supplies properness and compact norm-square critical set. The source's missing word-reference ``it'' prevents identifying a theorem about all toric varieties.
\subsection{Sources}
\begin{itemize}
\item Permanent UnsolvedMath ID \(20001976\), exact public code catalogue entry; local registry \texttt{aim-geometry-notes.json}, \texttt{notes[313]}, canonical tag question.
\item AIM workshop ``Moment maps and surjectivity in various geometries,'' Question 10.2, attributed to R. Goldin: \url{https://aimath.org/WWN/momentmaps/momentmaps.pdf}. No PDF was accessed; the exact OCR wording has no antecedent for ``it.''
\item The aid invokes a noncompact extension of Pomerleano--Teleman and the standard effective Delzant/Stanley--Reisner presentation, without full references or exact theorem citations. No original literature entry or year is supplied.
\item Research attempt 1, ``Split integral Kirwan surjectivity for smooth compact toric reductions,'' labelled partial\_result, partially\_solved, proved and low novelty confidence. The statement-verification field merely introduces the exact canonical OCR record.
\end{itemize}
\end{document}
